
\subsection{结果陈述}
\label{subsec:7-2-4}

现在实行第~7.2.1--7.2.3 小节提出的计划。
更具体地说，第~7.2.4--7.2.8 小节将完成
定理~\ref{thm:3-4-8} 与命题~\ref{prop:3-5-2} 的证明。
本小节精确陈述主要结果，
其证明放在下一小节。

将在 $L$ 上构造光滑奇异单形集
$\mathcal X_g(L)$（$g=0,1,2,\ldots$），
使其具有下列性质。记
\begin{equation}
\mathcal X_{(g_0)}(L)
=\bigcup_{g\le g_0}\mathcal X_g(L).
\tag{7.2.26}\label{eq:7-2-26}
\end{equation}
以 $\mathcal X_{(g)}(L)$ 为生成元的自由阿贝尔群
记为 $C_{(g)}(L;\mathbb Z)$，
以之为生成元的 $\mathbb Q$-向量空间
记为 $C_{(g)}(L;\mathbb Q)$。
令 $S(L;\mathbb Q)$ 为系数在 $\mathbb Q$ 中的
光滑奇异链复形。

\begin{statement}{性质}{7.2.27}\label{prop:7-2-27}
\liftIndex{I240}%
\noindent\textup{(7.2.28.1)}\quad
$\mathcal X_g(L)$ 是光滑奇异单形的可数集。

\noindent\textup{(7.2.28.2)}\quad
当 $g\ge1$ 时，包含映射
$C_{(g)}(L;\mathbb Q)\to S(L;\mathbb Q)$
诱导上同调同构。

\noindent\textup{(7.2.28.3)}\quad
$\mathcal X_g(L)$ 中每个元素 $P$ 的每个面
都属于 $\mathcal X_{(g)}(L)$。
\end{statement}

若 $P\in\mathcal X_g(L)$，
就称 $P$ 的\emph{代}\liftIndex{I244}为 $g$。
接着作如下定义。

\begin{definition}{7.2.29}\label{def:7-2-29}
设 $\beta\in\pi_2(M;L)$。
$\mathcal M^{\mathrm{main}}_{k+1}(\beta)$
的元素 $(\Sigma,w)$ 的一个\emph{装饰}\liftIndex{I243}，
是映射
$\mathfrak d\colon\{1,\ldots,k\}\to\mathbb Z_{\ge0}$。
\end{definition}

把 $\mathcal M^{\mathrm{main}}_{k+1}(\beta)$
的元素 $(\Sigma,w)$ 连同装饰 $\mathfrak d$
组成的集合记为
$\mathcal M^{\mathrm{main},\mathfrak d}_{k+1}(\beta)$。
由于其底层模空间与
$\mathcal M^{\mathrm{main}}_{k+1}(\beta)$ 相同，
它也带有仓西结构。
其元素写作 $(\Sigma,w,\mathfrak d)$。

置
\begin{equation}
\mathfrak G_0(L)=\{\beta\in\Pi(M;L)\mid
\mathcal M(\beta)\ne\varnothing\},
\tag{7.2.30}\label{eq:7-2-30}
\end{equation}
并令 $\mathfrak G(L)$ 为 $\mathfrak G_0(L)$ 在 $\Pi(M;L)$ 中
生成的子幺半群。
关于 $\Pi(M;L)=\Pi(L)$，见定义~\ref{def:2-4-17}。
记 $\beta_0=0\in\mathfrak G(L)$。
对 $\beta\ne\beta_0$、$\beta\in\mathfrak G(L)$，定义
\begin{equation}
\|\beta\|=
\sup\left\{n\ \middle|\
\begin{gathered}
\exists\,\beta_1,\ldots,\beta_n
\in\mathfrak G(L)\setminus\{\beta_0\},\\[-2pt]
\displaystyle\sum_{i=1}^{n}\beta_i=\beta
\end{gathered}\right\}
+\lfloor\omega(\beta)\rfloor-1,
\tag{7.2.31}\label{eq:7-2-31}
\end{equation}
其中 $\lfloor x\rfloor$ 是不超过 $x$ 的最大整数；
另置
\begin{equation}
\|\beta_0\|=-1.
\tag{7.2.32}\label{eq:7-2-32}
\end{equation}
条件 3.1.6 与式~\eqref{eq:7-2-31} 表明，
对任意给定的 $C$，满足 $\|\beta\|\le C$
的 $\beta\in\mathfrak G(L)$ 只有有限个。

对给定装饰
$\mathfrak d\colon\{1,\ldots,k\}\to\mathbb Z_{\ge0}$
及 $\beta\in\mathfrak G(L)$，若 $k\ne0$，定义
\begin{equation}
\|(\mathfrak d,\beta)\|
=\max_{i\in\{1,\ldots,k\}}\mathfrak d(i)
+\|\beta\|+k;
\tag{7.2.33}\label{eq:7-2-33}
\end{equation}
若 $k=0$，定义
\begin{equation}
\|(\mathfrak d,\beta)\|=\|\beta\|.
\tag{7.2.34}\label{eq:7-2-34}
\end{equation}
注意当 $\beta=\beta_0$ 时，
稳定性条件蕴含 $k\ge2$。

现在可以陈述本小节和下一小节的主要结果。
命题~\ref{prop:7-2-35} 中的
$\mathfrak s_{\mathfrak d,k,\beta,\vec P}$
在命题~\ref{prop:3-5-2} 中记为
$\mathfrak s_{\beta,\vec P}$。

\begin{proposition}{7.2.35}\label{prop:7-2-35}
对任意 $\delta>0$ 和 $K>0$，
存在 $\mathcal X_{(g)}(L)$（$g=0,\ldots,K$）
以及当 $\|(\mathfrak d,\beta)\|\le K$ 时的
多值截面 $\mathfrak s_{\mathfrak d,k,\beta,\vec P}$，
满足：

\noindent\textup{(7.2.36.1)}\quad
$\mathcal X_{(g)}(L)$ 满足性质~\ref{prop:7-2-27}。

\noindent\textup{(7.2.36.2)}\quad
设 $P_i\in\mathcal X_{\mathfrak d(i)}(L)$
（$i=1,\ldots,k$），置
\[
\mathcal M^{\mathrm{main},\mathfrak d}_{k+1}
(\beta;P_1,\ldots,P_k)
=\mathcal M^{\mathrm{main},\mathfrak d}_{k+1}(\beta)
\mathop{\times}_{L^k}\prod_{i=1}^{k}P_i.
\]
在其上定义多值截面
$\mathfrak s_{\mathfrak d,k,\beta,\vec P}$；
它与零截面横截。

\noindent\textup{(7.2.36.3)}\quad
若 $g=\|(\mathfrak d,\beta)\|$，则
\[
\operatorname{ev}_{0*}
\left(
\mathcal M^{\mathrm{main},\mathfrak d}_{k+1}
(L;\beta;P_1,\ldots,P_k)
^{\mathfrak s_{\mathfrak d,k,\beta,\vec P}}
\right)
\]
可分解为 $\mathcal X_{(g)}(L)$ 中的元素。
此处及下文用记号
\[
\mathcal M^{\mathrm{main},\mathfrak d}_{k+1}
(L;\beta;P_1,\ldots,P_k)
^{\mathfrak s_{\mathfrak d,k,\beta,\vec P}}
:=\mathfrak s_{\mathfrak d,k,\beta,\vec P}^{-1}(0)
\]
表示零集。

\noindent\textup{(7.2.36.4)}\quad
这些多值截面满足下述相容条件
7.2.38 与 7.2.44。
\translatornote{原书在此误作“7.2.38 and 30.44”；
本节后文第二条相容条件的可见编号为 7.2.44。
参见疑误 A-0435。}

\noindent\textup{(7.2.36.5)}\quad
$\mathfrak s_{\mathfrak d,k,\beta,\vec P}^{-1}(0)$
位于原仓西映射的零集的 $\delta$-邻域内。
\end{proposition}

严格说来，条件 (7.2.36.3)
需要以类似 (3.4.6.3) 的方式陈述：
零集 $\mathfrak s_{\mathfrak d,k,\beta,\vec P}^{-1}(0)$
具有一个三角剖分，
使其中每个单形都确定
$\mathcal X_{(g)}(L)$ 中的元素。
后文不再重复这一说明。

下面精确陈述多值截面的两条相容条件。
第一条描述其在奇异单形 $P_i$ 的边界附近的行为。
设 $P_i$ 由光滑映射
$f\colon|P_i|=\Delta^{d_i}\to L$ 给出；
对每个 $m$ 取 $p_{i,m}\in\Delta^{d_i}$，并设
\begin{equation}
\lim_{m\to\infty}p_{i,m}=p_{i,\infty}.
\tag{7.2.37}\label{eq:7-2-37}
\end{equation}
假设对某个 $i$，
$p_{i,\infty}\in\partial\Delta^{d_i}$。
进一步假设 $p_{i,\infty}$
在某个面 $\Delta^{d_i'}$ 的内部。
利用标准单形顶点的典范次序，
可将该面典范地识别为标准 $d_i'$-单形。
把 $(\Delta^{d_i'},f)$ 视为光滑奇异单形。
由 (7.2.28.3)，存在 $g'\le\mathfrak d(i)$，
使 $(\Delta^{d_i'},f)\in\mathcal X_{g'}(L)$。
记 $P_i'=(\Delta^{d_i'},f)$，
$\vec P'=(P_1',\ldots,P_k')$；
其中某些 $P_i'$ 可以仍等于 $P_i$。

现在说明上述极限下
$\mathfrak s_{\mathfrak d,k,\beta,\vec P}$ 的行为。
$\mathcal M^{\mathrm{main},\mathfrak d}_{k+1}(\beta)$
上的仓西结构按定义 A1.13 是弱浸没的，
因而满足以下性质。
设 $\mathbf p\in
\mathcal M^{\mathrm{main},\mathfrak d}_{k+1}(\beta)$，
$U_{\mathbf p}$ 为其仓西邻域；
又设
\[
(\mathbf p,(p_1,\ldots,p_k))
\in\mathcal M^{\mathrm{main},\mathfrak d}_{k+1}(\beta)
\mathop{\times}_{L^k}\prod_{i=1}^{k}P_i,
\]
且 $U_i(p_i)$ 是 $P_i$ 中 $p_i$ 的邻域。
则
$U_{\mathbf p}\mathop{\times}_{L^k}
\prod_iU_i(p_i)$
是 $(\mathbf p,(p_1,\ldots,p_k))$ 在
$\mathcal M^{\mathrm{main},\mathfrak d}_{k+1}
(\beta;P_1,\ldots,P_k)$ 中的仓西邻域。
设
$(\mathbf p_m,(p_{1,m},\ldots,p_{k,m}))$
位于该模空间的某个仓西邻域内，
且 $\lim_{m\to\infty}\mathbf p_m=\mathbf p_\infty$。

\begin{statement}{相容条件}{7.2.38}\label{cond:7-2-38}
\liftIndex{I241}%
设 $\mathbf p_m,\mathbf p_\infty$ 如上，
并取满足
$P_i'\in\mathcal X_{\mathfrak d'(i)}(L)$
的装饰 $\mathfrak d'$。则
\[
\lim_{m\to\infty}
\mathfrak s_{\mathfrak d,k,\beta,\vec P}
(\mathbf p_m,(p_{1,m},\ldots,p_{k,m}))
=\mathfrak s_{\mathfrak d',k,\beta,\vec P'}
(\mathbf p_\infty,(p_{1,\infty},\ldots,p_{k,\infty})).
\]
\end{statement}

第二条相容条件涉及
$\mathcal M^{\mathrm{main},\mathfrak d}_{k+1}(\beta)$
的边界。
设 $(\Sigma,w,\mathfrak d)$ 为其元素，
$\Sigma$ 奇异且边界上有二重点。
取其中一个二重点 $p_0\in\partial\Sigma$，
在此切开 $\Sigma$，得到 $\Sigma_1,\Sigma_2$；
不妨设第 $0$ 个标记点 $z_0$ 属于 $\Sigma_2$。
把 $w$ 在 $\Sigma_i$ 上的限制记作 $w_i$。
于是
$(\Sigma_i,w_i)\in
\mathcal M^{\mathrm{main}}_{k_i+1}(\beta_i)$，
其中 $\beta_1+\beta_2=\beta$、
$k_1+k_2-1=k$。记
\begin{equation}
(\Sigma,w)
=(\Sigma_1,w_1)\mathbin{\#_\ell}(\Sigma_2,w_2),
\tag{7.2.39}\label{eq:7-2-39}
\end{equation}
其中 $p_0$ 是 $\Sigma_2$ 的第 $\ell$ 个标记点，
也是 $\Sigma_1$ 的第 $0$ 个标记点。

\par\medskip\noindent\begin{minipage}{\textwidth}
\centering
\setlength{\unitlength}{0.8mm}
\begin{picture}(145,66)
\linethickness{0.5pt}
\put(50,31){\circle{42}}
\put(92,31){\circle{42}}
\linethickness{0.25pt}
\qbezier[16](50,57)(24,57)(24,31)
\qbezier[16](24,31)(24,5)(50,5)
\qbezier[10](50,5)(61,5)(68,12)
\qbezier[16](92,57)(118,57)(118,31)
\qbezier[16](118,31)(118,5)(92,5)
\qbezier[10](92,5)(81,5)(74,12)
\linethickness{0.5pt}
\put(71,31){\circle*{1.7}}
\put(50,31){\makebox(0,0){$\Sigma_1$}}
\put(92,31){\makebox(0,0){$\Sigma_2$}}
\put(73,58){$\scriptstyle p_0$}
\put(69,55){\line(1,-7){3}}
\put(80,48){\makebox(0,0){$\scriptstyle\times$}}
\put(79,51){$\scriptstyle z_{\ell-1}$}
\put(57,49){\makebox(0,0){$\scriptstyle\times$}}
\put(62,47){$\scriptstyle z_\ell$}
\put(87,51){\makebox(0,0){$\scriptstyle\times$}}
\put(104,45){\makebox(0,0){$\scriptstyle\times$}}
\put(107,45){$\scriptstyle z_1$}
\put(113,31){\makebox(0,0){$\scriptstyle\times$}}
\put(117,30){$\scriptstyle z_0$}
\put(108,17){\makebox(0,0){$\scriptstyle\times$}}
\put(111,14){$\scriptstyle z_k$}
\put(76,17){\makebox(0,0){$\scriptstyle\times$}}
\put(77,11){$\scriptstyle z_{\ell+k_1}$}
\put(58,13){\makebox(0,0){$\scriptstyle\times$}}
\put(61,7){$\scriptstyle z_{\ell+k_1-1}$}
\put(33,15){\makebox(0,0){$\scriptstyle\times$}}
\put(29,31){\makebox(0,0){$\scriptstyle\times$}}
\put(50,52){\makebox(0,0){$\scriptstyle\times$}}
\end{picture}
\par{\renewcommand{\thefigure}{7.2.1}\refstepcounter{figure}
\normalfont\bfseries 图~\thefigure\label{fig:7-2-1}}
\end{minipage}\par\medskip

接着确定 $(\Sigma_1,w_1)$、
$(\Sigma_2,w_2)$ 的装饰。
对 $i\in\{1,\ldots,k_1\}$，
设 $\Sigma_1$ 的第 $i$ 个标记点
对应于 $\Sigma$ 的第 $j$ 个标记点，置
\begin{equation}
\mathfrak d_1(i)=\mathfrak d(j).
\tag{7.2.40.1}\label{eq:7-2-40-1}
\end{equation}
再对 $i\in\{1,\ldots,k_2\}$，
若 $i\ne\ell$，设 $\Sigma_2$ 的第 $i$ 个标记点
对应于 $\Sigma$ 的第 $j$ 个标记点。
定义装饰 $\mathfrak d_2^{\max}$ 为
\begin{equation}
\mathfrak d_2^{\max}(i)=\mathfrak d(j)
\qquad(i\ne\ell),
\tag{7.2.40.2}\label{eq:7-2-40-2}
\end{equation}
\begin{equation}
\mathfrak d_2^{\max}(\ell)
=\|(\mathfrak d_1,\beta_1)\|.
\tag{7.2.40.3}\label{eq:7-2-40-3}
\end{equation}
这里 $\mathfrak d_1$ 已由 (7.2.40.1) 确定。
据此置
\begin{equation}
(\Sigma,w,\mathfrak d)
=(\Sigma_1,w_1,\mathfrak d_1)
\mathbin{\#_\ell}
(\Sigma_2,w_2,\mathfrak d_2^{\max}).
\tag{7.2.41}\label{eq:7-2-41}
\end{equation}
式~\eqref{eq:7-2-41} 给出映射
\begin{equation}
\begin{aligned}
\#_\ell\colon\quad
&\mathcal M^{\mathrm{main},\mathfrak d_1}_{k_1+1}(\beta_1)
\mathop{\times}_{\operatorname{ev}_0,\operatorname{ev}_\ell}
\mathcal M^{\mathrm{main},\mathfrak d_2^{\max}}_{k_2+1}(\beta_2)\\
&\longrightarrow
\partial\mathcal M^{\mathrm{main},\mathfrak d}_{k+1}(\beta).
\end{aligned}
\tag{7.2.42}\label{eq:7-2-42}
\end{equation}
命题~\ref{prop:7-1-2} 表明，
沿 $\#_\ell$ 把
$\mathcal M^{\mathrm{main},\mathfrak d}_{k+1}(\beta)$
的仓西结构限制到式~\eqref{eq:7-2-42} 左端的纤维积，
所得结构与该纤维积的仓西结构一致。

现在设
$P_i\in\mathcal X_{\mathfrak d(i)}(L)$
（$i=1,\ldots,k$）。考虑
\[
Q=\operatorname{ev}_{0*}\left(
\mathcal M^{\mathrm{main},\mathfrak d_1}_{k_1+1}
(\beta_1;P_\ell,\ldots,P_{\ell+k_1-1})
^{\mathfrak s_{\mathfrak d_1,k_1,\beta_1,\vec P_{(1)}}}
\right),
\]
其中
$\vec P_{(1)}=(P_\ell,\ldots,P_{\ell+k_1-1})$。
条件 (7.2.36.3) 蕴含，
$Q$ 可分解为
$\mathcal X_{(\mathfrak d_2^{\max}(\ell))}(L)$
中元素之和。
为使记号简洁，以下也用 $Q$
代表这一分解中出现的单个单形；
后文常用这种记号滥用，不再反复说明。
置
\begin{equation}
\mathfrak d_2(i)=
\begin{cases}
\mathfrak d_2^{\max}(i),&i\ne\ell,\\
g,&i=\ell,
\end{cases}
\tag{7.2.43}\label{eq:7-2-43}
\end{equation}
其中 $g$ 是 $Q$ 的代。再置
\[
\vec P_{(2)}
=(P_1,\ldots,P_{\ell-1},Q,
P_{\ell+k_1},\ldots,P_k),
\]
并考虑
$\mathfrak s_{\mathfrak d_2,k_2,\beta_2,\vec P_{(2)}}$。

注意
$\mathcal M^{\mathrm{main},\mathfrak d}_{k+1}(\beta)$
的仓西结构在其边界上的障碍丛，
可识别为
$\mathcal M^{\mathrm{main},\mathfrak d_1}_{k_1+1}(\beta_1)$
与
$\mathcal M^{\mathrm{main},\mathfrak d_2^{\max}}_{k_2+1}(\beta_2)$
的障碍丛
$E(\beta_1)_{k_1+1}$ 和 $E(\beta_2)_{k_2+1}$
的直和。
这由纤维积仓西结构的定义得出。
于是第二条相容条件可表述如下。

\begin{statement}{相容条件}{7.2.44}\label{cond:7-2-44}
\liftIndex{I242}%
经 $\#_\ell$ 拉回
$\mathfrak s_{\mathfrak d,k,\beta,\vec P}$，
把所得截面分解为
$s^{(1)}\oplus s^{(2)}
\in E(\beta_1)_{k_1+1}\oplus E(\beta_2)_{k_2+1}$。
则：

\noindent\textup{(1)}\quad
$s^{(1)}
=\mathfrak s_{\mathfrak d_1,k_1,\beta_1,\vec P_{(1)}}$。

\noindent\textup{(2)}\quad
在
$s^{(1)}
=\mathfrak s_{\mathfrak d_1,k_1,\beta_1,\vec P_{(1)}}$
为零的点，
$s^{(2)}
=\mathfrak s_{\mathfrak d_2,k_2,\beta_2,\vec P_{(2)}}$。
\end{statement}

$\mathfrak s_{\mathfrak d_2,k_2,\beta_2,\vec P_{(2)}}$
只在
$\mathfrak s_{\mathfrak d_1,k_1,\beta_1,\vec P_{(1)}}$
为零的点有定义。
换言之，它只在与
$\mathcal M^{\mathrm{main},\mathfrak d_2^{\max}}_{k_2+1}
(\beta_2;\vec P_{(2)})$
的仓西邻域中的点相对应之处有定义。
这里 $\vec P_{(2)}$ 包含 $Q$，
而 $Q$ 是
$\mathfrak s_{\mathfrak d_1,k_1,\beta_1,\vec P_{(1)}}$
的零集的虚基本链；
这就是条件 (2) 只作上述限定的原因。
在 $s^{(1)}$ 不为零的点，
$s^{(2)}$ 的值与
$\mathfrak s_{\mathfrak d,k,\beta,\vec P}$ 的零集无关，
因而这一限定对构造没有影响。

至此已精确陈述命题~\ref{prop:7-2-35}。
下一小节证明它。
